\section{The Container}
\label{sec:container}

This section describes the on-disk format the engine executes over. The format
and its compressor are prior work~\cite{gfaz2026}; we summarize only what the
engine relies on, in the terms of Section~\ref{sec:principles}, and we treat its
compression ratio and compressor throughput as given rather than as results.
Readers interested in how the grammar is constructed, in the GPU compression
backend, or in the encoding of the non-traversal record types should consult that
work.

\paragraph{Layout.}
The container is columnar (P3). Each GFA record type is split into typed columns,
which are transformed with a type-appropriate encoding and then entropy-coded with
Zstd~\cite{collet2021rfc} as independent blocks. A query decompresses only the
blocks it names. The columns that matter for this paper are the P and W traversal
payloads, their per-record encoded-length columns, the grammar rulebook, the
segment length column, the segment DNA payload, the link (edge) columns, and the
PanSN-style metadata columns that identify each traversal's sample and haplotype.

\paragraph{Traversal encoding.}
A traversal is a sequence of orientation-signed node IDs. Two transforms are
applied before entropy coding. The first is delta encoding, which exploits the
fact that consecutive steps in a pangenome walk usually land on nearby node IDs.
The second, which does the real work, is an iterative 2-mer grammar: repeated
adjacent symbol pairs are replaced by rule symbols over several rounds, so a
haplotype's walk collapses into a short string of rules that stand for long
shared substrings. Because human haplotypes agree over long stretches, the rules
are shared across samples, and this cross-record sharing is where the bulk of the
compression comes from.

The rulebook is stored as two parallel \texttt{int32} arrays, \texttt{first} and
\texttt{second}, giving each rule's two children. It is global to the file. The
encoded traversals are not: each is a self-contained string of literals and rule
references, delimited by the encoded-length column. This split is the design
decision Section~\ref{sec:principles} flagged as the expensive one. It preserves
record independence (P4) while still capturing cross-record redundancy, at the
cost of a rulebook that every query must load, and it is what allows an arbitrary
subset of haplotypes to be decoded by an arbitrary number of threads without
touching one another.

\paragraph{Namespace partitioning.}
Rule IDs are allocated sequentially from a \texttt{start\_id} that lies strictly
above every literal node ID in the file, and the layer metadata records the first
rule ID and the rule count. A decoder therefore classifies a symbol $e$ by testing
whether $|e|$ falls in the half-open rule range: in range means rule, out of range
means literal node ID (P2). No side table is consulted, and the test compiles to
two comparisons.

\paragraph{Orientation as sign.}
The sign of a symbol carries DNA strand. A negative literal is the segment read on
the reverse strand. A negative \emph{rule} reference denotes the reverse complement
of that rule: its children are visited in reverse order with their orientations
flipped, recursively. One rulebook therefore serves both strands, and no mirrored
copy is stored (P5).

\paragraph{What the engine assumes.}
The engine depends on four things and nothing else: that the rule range is
contiguous and disjoint from the literal range; that the rulebook is a static
array indexable by rule ID; that each traversal's encoded extent is recoverable
from a length column; and that rule negation means reverse complement. Any format
with those four properties can host the engine of Section~\ref{sec:engine}
unchanged. The engine does not depend on how the grammar was built, on the number
of grammar rounds, or on which entropy coder wraps the blocks.

\paragraph{Scale.}
For calibration, the container compresses the evaluation graphs by
$18$--$84\times$~\cite{gfaz2026}: the $369$\,GB HPRC v2.1 graph becomes
$4.5$\,GB, the $80$\,GB HGSVC3 graph $2.85$\,GB, and the $7.6$\,GB chr1 graph
$214$\,MB. These ratios are not claims of this paper; they set the scale of what
the engine reads.

\paragraph{An improved CPU backend.}
One change to the container is ours and is reported here for the first time. We
re-engineered the CPU compression and decompression paths, which now reach
\needsexp{X}\,GB/s and \needsexp{Y}\,GB/s respectively, against the
$385$\,MB/s and $1.8$\,GB/s reported for the same backend
in~\cite{gfaz2026}\needsexp{re-measure both configurations at a matched thread
count and under one methodology before printing any factor. The published
figures were taken at 16 threads with the page cache primed and averaged over
five runs; our current numbers are 32-thread end-to-end CLI wall times including
parse and serialization, so roughly half of the apparent gain is SMT rather than
optimization. Report the matched-thread factor, and report it per dataset:
whole-genome decompression improves far more than v1.1 compression does}.

We do not treat this as a contribution of the paper, and it does no work in the
results of Section~\ref{sec:eval}, where every \sys measurement is a query and
not a compression run. It matters here for one reason. The argument of
Section~\ref{subsec:bg-alternatives} is that expanding to fast local storage does
not rescue the baselines, and decompression throughput is what makes that
concrete: at these rates the container can be read from storage and
entropy-decoded faster than an NVMe device can stream the expanded GFA text,
so the expanded form is not cheaper to reach even when it is already on the
fastest tier available.
